RailFlux Interlocking Mechanisms - Technical Report
Signal-to-Signal Interlocking Architecture
Control Dependency Framework: The system implements hierarchical signal control through a JSON-based rule engine where signals establish controlling relationships. Each signal can be classified as independent (operates without dependencies) or controlled (requires permission from upstream signals). The InterlockingRuleEngine processes these relationships using two primary control modes:
AND Control Mode: All controlling signals must permit the requested aspect. If signal HM001 is controlled by signals ST001 and AS001 in AND mode, both controlling signals must show compatible aspects before HM001 can display a proceed indication.
OR Control Mode: Any single controlling signal can authorize the aspect change. This provides operational flexibility where multiple approach routes can independently authorize destination signal aspects.
Rule Evaluation Process: When a signal aspect change is requested, the system retrieves the current composite aspect of each controlling signal, evaluates conditional requirements (point machine positions, track clearances), and applies the control logic. Rules are structured as conditional permissions where controlling signal aspects combined with infrastructure states determine allowable subordinate signal aspects.
Composite Aspect Handling: The system processes both simple aspects (RED, YELLOW, GREEN) and composite aspects that combine main signal states with subsidiary indications (calling-on, loop signals). This enables complex signaling scenarios where auxiliary signal systems provide additional operational capabilities beyond basic three-aspect signaling.
Signal-to-Point Machine Interlocking
Conditional Aspect Authorization: Signal aspects are intrinsically linked to point machine positions through condition-based rules. A signal requesting GREEN aspect must verify that all associated point machines are in positions compatible with the intended train movement. The system evaluates point machine states as conditions within interlocking rules.
Position Dependency Validation: Before authorizing signal aspect changes, the InterlockingRuleEngine queries current point machine positions from the database and compares them against rule requirements. If a rule specifies point machine PM001 must be in NORMAL position for signal ST001 to show GREEN, the system blocks the aspect change if PM001 is in REVERSE position.
Dynamic State Synchronization: Point machine position changes trigger re-evaluation of dependent signal aspects. When operators change point positions, the system automatically assesses whether existing signal aspects remain valid under the new track configuration and can initiate protective aspect changes if conflicts arise.
Safety Interlocking: Point machine operations require protecting signals to be at RED aspect before movement is permitted. This prevents train movements over points during position changes and ensures mechanical safety of the switching infrastructure.
Signal-to-Circuit Protection Mechanisms
Protected Circuit Validation: Each signal maintains an array of protected track circuits that must be verified clear before proceed aspects can be displayed. The system implements dual-source validation, cross-referencing protected circuits defined in signal configuration data against those specified in interlocking rules to ensure data consistency.
Occupancy Detection Integration: Track circuit occupancy states are continuously monitored through database polling and real-time notifications. When a track circuit transitions from clear to occupied, the system immediately evaluates which signals protect that circuit and can automatically force those signals to RED aspect as a safety measure.
Multi-Circuit Protection: Signals can protect multiple track circuits simultaneously, enabling comprehensive route protection. The system validates all circuits in the protection array before authorizing proceed aspects, ensuring complete path clearance before trains are signaled to move.
Virtual Segment Abstraction: Track circuits aggregate into logical track segments that represent operationally meaningful sections of railway infrastructure. This abstraction layer allows signal protection logic to operate on logical track sections while maintaining connection to physical circuit hardware.
Signal Rendering and Visualization
Component Architecture: Signals are rendered as QML components with dynamic aspect visualization supporting the full spectrum of railway signal indications. The visual representation includes multiple aspect lights (typically three for main signals) with color states corresponding to operational aspects.
Aspect State Mapping: The rendering system maps database aspect values to visual representations, translating text-based aspect codes (RED, YELLOW, GREEN, WHITE, BLUE) into appropriate color displays and lighting patterns on the signal head visualization.
Real-Time Updates: Signal visual states update immediately upon database state changes through Qt's signal-slot mechanism. Database notifications trigger QML property updates that drive visual aspect changes without requiring explicit refresh operations.
Signal Type Differentiation: Different signal types (OUTER, HOME, STARTER, ADVANCED_STARTER) are rendered with appropriate visual characteristics reflecting their operational roles and aspect capabilities. Advanced signals may display additional subsidiary aspects beyond the basic three-aspect configuration.
Circuit-Segment Interaction Layer
Virtual Abstraction Layer: Track segments represent logical groupings of physical track circuits, providing operational abstractions that align with railway operational practices. A single track segment may encompass multiple track circuits, aggregating their states into unified occupancy information.
State Aggregation Logic: Track segment occupancy is determined through aggregation of constituent circuit states. If any circuit within a segment reports occupied status, the entire segment is considered occupied. This provides operators with simplified infrastructure state information while maintaining detailed circuit-level monitoring for safety systems.
Protection Relationship Mapping: Track segments maintain relationships to protecting signals, establishing which signals provide safety coverage for each logical track section. This enables both forward and reverse lookup operations for safety validation and automatic protection activation.
Dynamic State Management: Segment states update automatically as constituent circuit states change. The system maintains consistency between physical circuit states and logical segment representations, ensuring operational accuracy in both detailed technical views and simplified operational displays.
Point Machine Visualization and Control
Position State Representation: Point machines are rendered as QML components displaying current position (NORMAL/REVERSE) through visual indicators that clearly distinguish between the two operational positions. The visualization includes directional indicators showing track connectivity under each position state.
Operational Status Display: Point machine components show operational status including normal operation, transition states, locked conditions, and failure modes. Visual indicators provide operators with immediate status information for infrastructure management.
Paired Machine Coordination: The system handles paired point machine operations where multiple machines must move simultaneously to maintain track geometry. Visual representations indicate paired relationships and coordinate position changes to ensure synchronized operation.
Safety State Visualization: Point machine components display safety-relevant information including protecting signal states, track circuit occupancy affecting movement authorization, and lock conditions preventing unauthorized operation. This provides operators with comprehensive infrastructure safety status at a glance.
Interactive Control Interface: Point machine components support operator interaction for position changes, incorporating interlocking validation feedback. The interface prevents unauthorized operations and provides immediate feedback on operation success or failure through visual state changes and status messaging.